\section{Related Work}
\label{sec:related}

\para{Protein language models and mutation-effect benchmarks.} Protein language models learn representations from large protein-sequence corpora and can support downstream biological prediction without direct supervision for each task. Earlier work showed that large protein sequence models encode biological structure and function \cite{rives2021biological}, and \esm-family models scale this idea to atomic-level structure prediction \cite{lin2023evolutionary}. \proteingym standardizes deep mutational scanning and clinical variant benchmarks for evaluating protein fitness and design models \cite{notin2023proteingym}. We use \proteingym because it provides direct mutation-effect measurements, making it suitable for asking whether a residue-level internal feature marks positions where mutations are experimentally deleterious.

\para{Sparse features in protein language models.} Sparse autoencoders have become a central tool for decomposing dense model activations into higher-dimensional sparse features. \interplm trains sparse autoencoders on \esm embeddings and shows that sparse features align with many curated concepts, while raw neurons show less conceptual alignment \cite{simon2025interplm}. \citet{gujral2025sparse} similarly use sparse autoencoders and transcoders to extract interpretable protein-level and amino-acid-level features from PLM representations. \citet{adams2025mechanistic} frame this line of work as a path from mechanistic interpretability to biological hypothesis generation. We build on this view but add a stricter empirical standard: an unannotated feature must be compared to raw hidden dimensions, shuffled sparse controls, and motif scans before it is treated as a candidate.

\para{Motif discovery and annotation controls.} Motifs provide a natural residue-level bridge between model internals and known biology. The \elm resource catalogs eukaryotic short linear motifs \cite{kumar2024elm}, and \prosite provides curated protein patterns, profiles, and rules for protein-domain and functional-site annotation \cite{sigrist2013prosite}. MotifAE adapts sparse autoencoders for motif discovery by encouraging local activation coherence and reports improved recovery of known \elm motifs \cite{hou2025motifae}. Our experiment uses \elm and \prosite differently: not as the prediction target, but as a conservative exclusion filter for claims of candidate novelty. We also scan a \swissprot sample from UniProtKB \cite{uniprot2025} to estimate whether candidate feature activations align with motif-positive residues outside the target assays.

\para{Why causal checks matter.} Feature-label association is not enough to infer mechanism. A feature can correlate with \dms sensitivity because of conservation, local structure, motif-database incompleteness, or the large number of tested features. Interventions provide a stronger test by asking whether changing the feature changes model computation. Our ablations are still internal model interventions, not wet-lab validation, but they help separate features that merely track a label from features that participate in \esm masked-marginal scoring.
